\section{Related Work}
\label{sec:relatedwork}

%Having presented our routing method, we situate this work within the broader research landscape from three perspectives: existing 

This work builds on works from three perspectives: filtered ANN methods (Section~\ref{sec:rw-fanns}), adaptive query planning and method selection (Section~\ref{sec:rw-routing}), and related benchmarks (Section~\ref{sec:rw-bench}).

\subsection{Filtered ANN Search}
\label{sec:rw-fanns}

Existing categorical-label Filtered ANN (FANNS) methods can be categorized by execution strategy into three classes~\cite{FilteredANNBenchmark2024}.

\textit{Filter-then-search} methods first apply the attribute constraint to retrieve a candidate subset, then perform ANN search on this subset. Representative methods include Pre-filter, which performs brute-force search on the filtered subset, and UNG~\cite{UNG2024}, which builds per-attribute sub-graph indices connected by a label-navigating graph for cross-partition traversal. These methods perform well under high selectivity (small candidate subsets) but degrade toward linear scan under low selectivity.

\textit{Search-then-filter} methods first perform ANN search on a global index, then apply the attribute constraint to the retrieved candidates. Representative methods include Post-filter HNSW~\cite{HNSW2018} and Post-filter IVFPQ~\cite{FaissGPUIVFPQ2019}. While they directly reuse mature ANN index structures, under strong filter selectivity few of the initial top-$k'$ candidates satisfy the constraint, requiring substantial candidate-set expansion to maintain recall, with significant efficiency degradation.

\textit{Hybrid-search} methods embed attribute constraints directly into ANN index construction or search. Early work HQANN~\cite{HQANN2022} introduced lightweight attribute-aware filtering on HNSW; Filtered-DiskANN and Stitched-DiskANN~\cite{FilteredStitchedVamana2023} introduce label-aware neighbor selection on the Vamana graph; ACORN-1 and ACORN-$\gamma$~\cite{ACORN2024} extend HNSW with attribute-aware pruning and multi-hop neighbor expansion; SIEVE~\cite{SIEVE2025} uses historical query workloads to offline-construct a heterogeneous collection of sub-indices targeting common filter patterns; CAPS~\cite{CAPS2023} and NHQ~\cite{NHQNPG2023} primarily target fixed-length-label scenarios, where CAPS partitions the attribute space via K-means and AFT, while NHQ fuses vector and attribute distances into a unified weighted distance. While these methods generally achieve balanced performance across filter scenarios, each retains structural preferences for specific filter types (containment, overlap, equality) and data distributions.

Beyond the categorical-label setting above, a parallel line of work targets \emph{range filters} over numerical attributes such as timestamps and prices, where a query consists of a vector $v_q$ and a continuous interval $[\ell, r]$. Representative methods include iRangeGraph~\cite{iRangeGraph2024} (range-dedicated graphs assembled from sub-graphs), SeRF~\cite{SeRF2024} (segment-graph augmentation of HNSW), ARKGraph~\cite{ARKGraph2023} (all-range $k$-NN graph), UNIFY~\cite{UNIFY2025} (unified index spanning the full range axis), WindowFilters~\cite{WindowFilters2024} (window-based precomputed indices), RangePQ~\cite{RangePQ2025} (efficient dynamic indexing for range queries), and TimestampANN~\cite{TimestampANN2025} (timestamp-specialized variant). As noted in the FANNS benchmark~\cite{FilteredANNBenchmark2024}, these methods are largely incompatible with the categorical-label algorithms above and are not the focus of this paper.

In industry, vector databases such as Milvus~\cite{Milvus2021} and PASE~\cite{PASE2020} have integrated FANNS as a core query capability, but their internal execution still relies on a single fixed strategy (e.g., Pre-filter or Post-filter), lacking adaptive cross-method selection.

A unified-benchmark cross-method evaluation~\cite{FilteredANNBenchmark2024} on categorical setting reveals that \emph{no single method dominates across all scenarios}: different methods alternately lead under different (dataset, query scenario) combinations, exhibiting strong structural complementarity. This observation directly motivates our investigation of adaptive method selection.

\subsection{Adaptive Query Planning and Method Selection}
\label{sec:rw-routing}

Our routing work fundamentally instantiates the broader \emph{algorithm selection} problem~\cite{Rice1976} in the FANNS domain. This paradigm has mature portfolio-based solutions in domains such as SAT solving~\cite{SATzilla2008}, where the optimal solver is predicted from problem-instance features rather than relying on a single universal method. Recently, similar ideas have begun extending to vector retrieval.

The most closely related work is the learning-based query planner of Gan and Wang~\cite{LearningBasedQueryPlanning2026}. Their method trains a two-layer MLP classifier per dataset to select between Pre-filter and Post-filter on a per-query basis, using lightweight features (selectivity, dimensionality, dataset distribution). On ArXiv and Wolt datasets, they achieve up to $4\times$ speedup over single-strategy baselines. However, their strategy space is limited to two naive methods (Pre-filter vs.\ Post-filter), excluding hybrid-search methods such as UNG, ACORN, and SIEVE; their planner is also trained per dataset, with no empirical evidence on cross-dataset generalization.

A complementary line of work comes from unfiltered VSS. Iceberg~\cite{Iceberg} addresses method selection for general vector similarity search from a task-centric view, proposing the Information Loss Funnel model and deriving an interpretable decision tree from easy-to-compute meta-features such as clustering tightness (DBI), vector norm coefficient of variation (CV), relative angle (RA), and relative contrast (RC), to guide selection among methods like HNSW, NSG, and RaBitQ. While Iceberg does not address attribute filtering, its methodology of \emph{constructing an interpretable decision tree from dataset meta-features} is structurally analogous to ours, validating meta-feature routing as a viable direction in the broader VSS area.

Our work differs from the above in three notable ways: (1) \emph{broader method coverage}, with the candidate pool spanning all three FANNS execution paradigms (filter-then-search, search-then-filter, hybrid-search) rather than only Pre vs.\ Post; (2) \emph{both rule-based and ML routers are provided}, with the former derived directly from structural analysis through a three-feature decision tree (query scenario, label cardinality, LID), and the latter approximating the Oracle upper bound via per-method regression modeling; (3) \emph{systematic evaluation of generalization on 5 mid-scale out-of-sample validation datasets} (500K--800K vectors, including yahoo and dbpedia real-text data), directly testing routing transferability beyond the training distribution.

\subsection{Benchmarks for Filtered ANN}
\label{sec:rw-bench}

The most directly related benchmark is the unified FANNS benchmark of Shi et al.~\cite{FilteredANNBenchmark2024}, which systematically evaluates 10 FANNS algorithms (NHQ, Filtered-DiskANN, Stitched-DiskANN, ACORN-1, ACORN-$\gamma$, CAPS, UNG, Pre-filter Brute-Force, Post-filter HNSW, Post-filter IVFPQ) on 6 real-world datasets (including YFCC and YouTube) under five filter modes (containment, overlap, equality, fixed-equality, combined). They categorize methods into the filter-then-search, search-then-filter, and hybrid-search paradigms, and emphasize parameter-fairness in evaluation. Our empirical foundation builds directly on this benchmark: on top of its method and dataset pool, we introduce dataset difficulty features (LID, RC, distribution factor) as routing signals and independently construct 5 mid-scale validation datasets (including yahoo and dbpedia real-text data) to assess routing generalization beyond the training set.

General ANN benchmarks such as ANN-Benchmarks~\cite{ANNBenchmarks2020} and Big-ANN-Benchmarks~\cite{simhadri2022bigann} are mature in the unfiltered setting, but their evaluation protocols do not cover attribute-filtered scenarios. Iceberg~\cite{Iceberg} extends end-to-end task-centric evaluation in the general VSS setting, but likewise does not address filter constraints. Our work complements this benchmark ecosystem by transitioning from benchmarking toward \emph{method selection grounded in structural observation} in the FANNS setting.